\documentclass[11pt,a4paper]{article}

\usepackage[utf8]{inputenc}
\usepackage[margin=1in]{geometry}
\usepackage{titlesec}
\usepackage{hyperref}
\usepackage{listings}
\usepackage{xcolor}
\usepackage{booktabs}
\usepackage{microtype}
\usepackage{fancyhdr}

% Page styling
\pagestyle{fancy}
\fancyhf{}
\rhead{\textbf{Arch Linux Peer-to-Peer Transfer}}
\lhead{Technical Documentation}
\rfoot{Page \thepage}

% Color palette for code & aesthetics
\definecolor{codegreen}{rgb}{0.2,0.6,0.2}
\definecolor{codegray}{rgb}{0.5,0.5,0.5}
\definecolor{codepurple}{rgb}{0.58,0,0.82}
\definecolor{codeblue}{rgb}{0.13,0.34,0.68}
\definecolor{backcolour}{rgb}{0.96,0.97,0.98}

\lstdefinestyle{codeblock}{
    backgroundcolor=\color{backcolour},   
    commentstyle=\color{codegreen},
    keywordstyle=\color{codeblue}\bfseries,
    numberstyle=\tiny\color{codegray},
    stringstyle=\color{codepurple},
    basicstyle=\ttfamily\footnotesize,
    breakatwhitespace=false,         
    breaklines=true,                 
    captionpos=b,                    
    keepspaces=true,                 
    numbers=left,                    
    numbersep=5pt,                  
    showspaces=false,                
    showstringspaces=false,
    showtabs=false,                  
    tabsize=2,
    frame=single,
    rulecolor=\color{codegray}
}

\lstdefinestyle{bashblock}{
    backgroundcolor=\color{backcolour},
    basicstyle=\ttfamily\small,
    breaklines=true,
    frame=single,
    rulecolor=\color{codegray}
}

\lstset{style=codeblock}

\hypersetup{
    colorlinks=true,
    linkcolor=codeblue,
    urlcolor=codeblue
}

\title{\textbf{\Huge Arch Linux Local Peer-to-Peer Folder \& File Transfer}\\[0.3em]
\Large Technical Reference, Architecture, and User Manual}
\author{\textbf{Bare-Minimum C++ Systems Implementation}}
\date{\today}

\begin{document}

\maketitle
\tableofcontents
\vspace{1.5em}
\hrule
\vspace{1.5em}

\section{Introduction and Overview}
Transferring files and directory trees between adjacent Arch Linux workstations on a Local Area Network (LAN) often involves friction—requiring flash drives, cloud storage intermediaries, or complex network share configurations (such as Samba or NFS).

This software provides a **bare-minimum, zero-external-dependency** tool implemented strictly in standard C++23, POSIX system calls, and minimal X11/XDnD. It allows two machines on the same WiFi or Ethernet network to:
\begin{enumerate}
    \item \textbf{Auto-discover} each other without manual IP configuration.
    \item \textbf{Switch roles} via a single button (acting as a background \textit{Client/Receiver} or an interactive \textit{Sender}).
    \item \textbf{Preserve full metadata}: File permissions (\texttt{chmod}), modification timestamps (\texttt{mtime}), and access timestamps (\texttt{atime}) are mirrored on the destination.
    \item \textbf{Stream recursively}: Entire nested folder trees and arbitrary files stream in chunked batches over raw TCP sockets.
    \item \textbf{Drag-and-Drop}: Drag files or directories directly from file managers (e.g., KDE Dolphin) into the application window.
\end{enumerate}

\section{Technology Stack \& Architecture}

\subsection{Architectural Principles}
\begin{itemize}
    \item \textbf{Language Standard}: C++23 (\texttt{g++ -std=c++23}).
    \item \textbf{No External Frameworks}: Zero third-party packages (no Qt, GTK, Boost, Asio, or Curl).
    \item \textbf{Display System}: Raw \texttt{libX11} client integration, natively compatible with Wayland (KDE Plasma) via XWayland.
    \item \textbf{Networking}: Raw Berkeley / POSIX socket interfaces (\texttt{AF\_INET}, \texttt{SOCK\_DGRAM} for discovery, \texttt{SOCK\_STREAM} for transfers).
\end{itemize}

\subsection{Communication Protocols}

\subsubsection{1. Automatic Discovery Subsystem (UDP)}
Nodes periodically broadcast heartbeat packets on the subnet broadcast address (\texttt{255.255.255.255}) over port \texttt{UDP 50000}:
\begin{itemize}
    \item \textbf{Interval}: Heartbeats fire every $2$ seconds.
    \item \textbf{Payload Format}: Packed binary structure containing protocol magic (\texttt{0x44495343}), listening TCP transfer port, role status (Client vs Sender), and local hostname.
    \item \textbf{Liveness Evaluation}: Active peers are refreshed in the list; peers silent for $>8$ seconds are pruned automatically.
\end{itemize}

\subsubsection{2. Multi-Item Stream Transfer Subsystem (TCP)}
Transfers execute over direct point-to-point connections on port \texttt{TCP 50001}:
\begin{itemize}
    \item Every file/directory is preceded by a packed \texttt{PacketHeader}:
    \begin{lstlisting}[language=C++]
struct PacketHeader {
    uint8_t entry_type;  // REGULAR_FILE = 1, DIRECTORY = 2, END = 255
    uint32_t path_len;   // UTF-8 path length
    uint64_t file_size;  // Size in bytes (0 for directory)
    uint32_t mode;       // st_mode permissions & file type
    int64_t  mtime_sec;  // Modification timestamp (sec)
    int64_t  mtime_nsec; // Modification timestamp (nsec)
    int64_t  atime_sec;  // Access timestamp (sec)
    int64_t  atime_nsec; // Access timestamp (nsec)
};
    \end{lstlisting}
    \item \textbf{Data Streaming}: Files are streamed in 64\,KB chunks (\texttt{65536} bytes), keeping memory usage constant regardless of file size.
    \item \textbf{Metadata Restoration}: Upon receiving the terminal \texttt{END\_OF\_BATCH} packet, the receiver applies filesystem timestamps (\texttt{utimensat}) and permission masks (\texttt{chmod}) in reverse topological order so directory attributes are preserved.
\end{itemize}

\subsubsection{3. Drag-and-Drop Engine (X11 / XDnD)}
Drag-and-drop utilizes the standard X11 Drag-and-Drop protocol version 5 (\texttt{XdndAware}, \texttt{XdndEnter}, \texttt{XdndPosition}, \texttt{XdndDrop}, \texttt{XdndSelection}):
\begin{itemize}
    \item Decodes dropped MIME payloads of type \texttt{text/uri-list}.
    \item Converts \texttt{file://} encoded URIs into native POSIX filesystem paths.
\end{itemize}

\section{Building from Source}

The program does not require system-wide installation. Compile with the included \texttt{Makefile}:

\begin{lstlisting}[style=bashblock]
# Navigate to the project directory
cd folder-sharing

# Build the standalone binary
make clean && make
\end{lstlisting}

This produces an optimized binary named \texttt{folder\_share}.

\section{Usage and Command Reference}

\subsection{Graphical Mode (GUI)}
To start in graphical mode (default):
\begin{lstlisting}[style=bashblock]
./folder_share
\end{lstlisting}

\subsubsection{Interface Modes}
\begin{itemize}
    \item \textbf{Client / Receiver Mode}: 
    \begin{itemize}
        \item Click the \texttt{[ Client / Receiver ]} button.
        \item The application listens on port \texttt{50001}. Any incoming file or directory tree is saved automatically to \texttt{\textasciitilde/Downloads}.
    \end{itemize}
    \item \textbf{Sender Mode}:
    \begin{itemize}
        \item Click the \texttt{[ Sender ]} button.
        \item The left pane lists all automatically discovered machines on the network.
        \item Click a target peer in the list.
        \item Drag and drop any file or folder from your file manager (Dolphin) into the dropzone on the right.
        \item The transfer starts immediately with real-time progress.
    \end{itemize}
\end{itemize}

\subsection{Command-Line Interface (CLI)}

The binary includes full CLI functionality for headless systems or terminal scripts.

\begin{table}[h!]
\centering
\small
\begin{tabular}{@{}ll@{}}
\toprule
\textbf{Command} & \textbf{Description} \\ \midrule
\texttt{./folder\_share --cli receive} & Runs as a receiver daemon; saves files to \texttt{\textasciitilde/Downloads}. \\
\texttt{./folder\_share --cli list} & Broadcasts and prints discovered network peers. \\
\texttt{./folder\_share --cli send <peer-idx|ip> <paths...>} & Sends files or folders to a discovered peer or IP. \\
\texttt{./folder\_share --help} & Displays command usage summary. \\ \bottomrule
\end{tabular}
\caption{CLI Command Reference}
\end{table}

\section{Step-by-Step Practical Use Cases}

\subsection{Use Case 1: GUI Drag-and-Drop Sharing}
\textbf{Scenario}: Desktop $A$ (\textit{Arch-Desktop}) wants to share a project directory \texttt{\textasciitilde/Projects/MyEngine} to Laptop $B$ (\textit{Arch-Laptop}).

\begin{enumerate}
    \item On Laptop $B$, launch the app:
    \begin{lstlisting}[style=bashblock]
./folder_share
    \end{lstlisting}
    Leave it on the default \texttt{[ Client / Receiver ]} mode.
    \item On Desktop $A$, launch the app:
    \begin{lstlisting}[style=bashblock]
./folder_share
    \end{lstlisting}
    Click \texttt{[ Sender ]}.
    \item Desktop $A$ immediately sees \texttt{Arch-Laptop (192.168.1.42:50001)} in the peer list. Click on it.
    \item Open Dolphin, drag the \texttt{MyEngine} folder, and drop it into the dropzone.
    \item The entire folder tree is transferred into \texttt{\textasciitilde/Downloads/MyEngine} on Laptop $B$ with permissions and timestamps intact.
\end{enumerate}

\subsection{Use Case 2: Terminal / Headless CLI Transfer}
\textbf{Scenario}: You are SSH'd into a remote Arch workstation and want to receive files from your local terminal.

\begin{enumerate}
    \item On the receiving machine:
    \begin{lstlisting}[style=bashblock]
./folder_share --cli receive
    \end{lstlisting}
    \textit{Output:}
    \begin{lstlisting}[style=bashblock]
[Client Mode] Starting receiver...
[Client Mode] Saving incoming files to: /home/niraj/Downloads
[Client Mode] Hostname: Arch-Receiver
[Client Mode] Listening on TCP port 50001 and broadcasting presence on UDP 50000
    \end{lstlisting}

    \item On the sender machine, list available peers:
    \begin{lstlisting}[style=bashblock]
./folder_share --cli list
    \end{lstlisting}
    \textit{Output:}
    \begin{lstlisting}[style=bashblock]
Discovered Peers (1 found):
--------------------------------------------------------
Index Hostname                  IP Address        Port    
--------------------------------------------------------
[1]   Arch-Receiver             192.168.1.42      50001   
--------------------------------------------------------
    \end{lstlisting}

    \item Send the directory using the discovered peer index \texttt{1}:
    \begin{lstlisting}[style=bashblock]
./folder_share --cli send 1 ~/Documents/Research ~/Pictures/photo.png
    \end{lstlisting}
    \textit{Output:}
    \begin{lstlisting}[style=bashblock]
[Sender] Connecting to 192.168.1.42:50001...
[↑] Sending: Research/report.pdf [4.20 MB / 4.20 MB] (100.0%)
[✓] Transfer complete! (4.20 MB sent)
    \end{lstlisting}
\end{enumerate}

\section{Security \& Firewall Considerations}
\begin{itemize}
    \item \textbf{Local Network Only}: By design, transfers operate over plaintext TCP for minimal overhead. Only use within trusted local area networks.
    \item \textbf{Firewall Ports}: If running \texttt{ufw} or \texttt{iptables}, ensure traffic is allowed on:
    \begin{itemize}
        \item \texttt{UDP 50000} (Inbound/Outbound LAN broadcast discovery).
        \item \texttt{TCP 50001} (Inbound file transfer connections on client).
    \end{itemize}
\end{itemize}

\end{document}
